% [原文 p.53]
\chapter{簇代数簇}

\ChapterAbstract{本章引入簇代数簇（cluster variety）与簇代数（cluster algebra）。簇代数簇容许一种特殊的图册（atlas），其中转移函数是某些称为变异（mutation）的双有理变换。许多与纤毛曲面（ciliated surface）相伴的模空间都具有这种结构，例如射影直线上 $n$ 个点的模空间。在与纤毛曲面相伴的某些空间上存在自然的坐标系，曲面的每个三角剖分对应其中之一。在这个意义上，一个三角剖分给出这类空间上的一个坐标卡（chart）；这里的「坐标卡」应按代数几何而非微分几何的意义来理解，因为每个坐标卡都覆盖空间的一个开稠密子集。簇代数提供了一种深刻的代数结构，它帮助我们理解两个这样的坐标卡之间的转移函数，或者更明确地说，理解如何把对应于同一个底纤毛曲面的两个不同三角剖分的坐标系联系起来。簇代数是相当晚近的代数结构，由 Fomin、Zelevinsky 与 Berenstein 在一系列论文 [133, 134, 48, 135] 中引入。}

\section{起始例子}

在引入簇代数簇的主要概念之前，让我们先看一个具体例子，它能很好地展示簇代数簇的整体结构。

我们考虑 $\C^2$ 中 $n\ge 3$ 条向量直线（vectorial line）模去 $\PGL_2(\C)$ 的自然作用，即 $\C\PP^1$ 上 $n$ 个点模去射影变换。把这个模空间记作 $\mathcal{M}_n$。

先考虑 $n=3$ 的情形。由于 $\PGL_2(\C)$ 在一般的三元直线组上是 3-可迁的，只要 $\ell_i\ne\ell_j$（对所有 $i\ne j$），$(\ell_1,\ell_2,\ell_3)$ 就只有一个轨道。当某些直线重合时还有 4 个其他轨道：三个恰有两条直线重合的轨道，以及 $(\ell_1,\ell_1,\ell_1)$。利用商拓扑可以注意到，这些非一般轨道的任何邻域都与一般轨道相交。这意味着朴素的拓扑商给出的空间不是豪斯多夫的（在我们这里：5 个点，其中 4 个与一般点无限接近）。

% [原文 p.54]
为解决这个问题，我们使用所谓的 GIT 商，以得到豪斯多夫的模空间 $\mathcal{M}_n$。

代数簇 $X$ 被作用在其上的群 $G$ 的 GIT 商是一个豪斯多夫空间 $X//G$，它具有如下泛性质：所有从 $X$ 到 $Y$（一个带 $G$ 作用的代数簇）的 $G$-等变映射都通过 $X//G$ 分解。实践中，GIT 商就是合并掉非豪斯多夫点的朴素拓扑商。对 $\mathcal{M}_3$ 而言，我们把两个非一般点与一般点合并。于是 $\mathcal{M}_3$ 只是一个点。

现在考虑 $n=4$。对点的四元组 $(P_1,P_2,P_3,P_4)$ 存在一个射影不变量，即交比
\[
(P_1,P_2,P_3,P_4):=\frac{P_1-P_3}{P_2-P_3}\cdot\frac{P_2-P_4}{P_1-P_4}\in\C\cup\{\infty\}
\]
（我们把 $\C\PP^1$ 与 $\C\cup\{\infty\}$ 等同）。对一般的四元组，即当 $P_i\ne P_j$（对所有 $i\ne j$）时，交比取值于 $\C\PP^1\setminus\{0,1,\infty\}$。其余的值 $0$、$1$ 与 $\infty$ 在 4 条直线中有两条重合时取到，这有 $\binom{4}{2}=6$ 种可能方式。每个特殊值恰好由两对取到。例如 $P_1=P_2$ 给出交比 $1$，而 $P_3=P_4$ 也给出交比 $1$。两对 $(P_1,P_1,P_3,P_4)$ 与 $(P_1,P_2,P_3,P_3)$ 无限接近（两个轨道的任意两个邻域都相交），因此它们在 GIT 商下被等同。最后，只要多于 2 个点重合，该轨道就与所有其他点无限接近，于是它们被其他点吸收。

最终结果是 $\mathcal{M}_4\cong\C\PP^1$，它有两个坐标卡，分别由交比 $x=(P_1,P_2,P_3,P_4)\in\C$ 与 $y=(P_2,P_3,P_4,P_1)\in\C$ 给出，二者通过变换 $y=x^{-1}$ 相联系。交比给出球面 $\C\PP^1\cong S^2$ 由北极与南极给出的通常坐标卡。这是簇代数簇的典型特征：每个坐标卡都开且稠密，而转移函数是特殊的有理映射。

% [Tutor Note: 交比约定与 $y=x^{-1}$ 的核对。原文取 $(P_1,P_2,P_3,P_4)=\frac{(P_1-P_3)(P_2-P_4)}{(P_2-P_3)(P_1-P_4)}$，这正是标准交比 $(P_1,P_2;P_3,P_4)$。取 $P_1=0,P_2=1,P_3=2,P_4=3$，得 $x=4/3$；而此时 $y=(P_2,P_3,P_4,P_1)=\frac{(P_2-P_4)(P_3-P_1)}{(P_3-P_4)(P_2-P_1)}=4=x/(x-1)$，并非 $x^{-1}$。若要 $y=x^{-1}$，应取 $y=(P_1,P_2,P_4,P_3)$（同一组 $P$ 给出 $3/4$）。两种取法只改变两坐标卡之间双有理过渡函数的表达式；后文（例 5.13 与式 (5.1)）只用其双有理性，故不影响结论，读者核对时以自己算出的过渡函数为准。]
为了看清纤毛曲面是如何出现的，最后让我们考虑 $n=5$。于是我们感兴趣的是 $\C^2$ 中 5 条向量直线的模空间。

考虑一个五边形，把其中一条直线对应到它的每个顶点，比如顶点 $i$ 对应直线 $\ell_i$。注意五边形是纤毛曲面的一个非常特殊的情形。然后固定五边形的一个三角剖分。按下述方式给两条对角线各配一个复数：一条对角线恰好位于一个四边形内。与这条对角线相伴的数是该四边形四条直线的交比（按顺时针方向，并从对角线上的一个点出发）。例如在图 5.1 中，$x=(\ell_1,\ell_3,\ell_4,\ell_5)$，$y=(\ell_1,\ell_2,\ell_3,\ell_4)$。交比的不变性表明这个数与起点无关。

我们断言这两个数构成 $\mathcal{M}_5$ 的一个坐标系。确实，我们可以由这两个数重构出这五条直线：由于 $\PGL_2(\C)$ 在 $\C^2$ 的直线集合上是 3-可迁的，我们可以固定第一个四边形中四条直线里的三条（比如 $\ell_1,\ell_2,\ell_3$）。第四条直线随即由写在对角线上的交比确定。第五条直线再由另一个交比确定。

我们看到，这一重构只给出一般的五元组直线。因此，由五边形的一个固定三角剖分给出的一个坐标卡，描述的是 $\mathcal{M}_5$ 的一个开稠密子集。

% [原文 p.55]
\Fig[0.5\linewidth]{figures/fig_5_1.png}{5.1}{五边形的一个固定三角剖分给出 $\mathcal{M}_5$ 上的坐标。}

当我们通过一次翻转（flip）改变三角剖分时，可以利用交比的性质确定坐标如何变化。翻转与 $x$ 相伴的对角线后，新坐标由 $(x^{-1},y(1+x))$ 给出。若翻转与 $y$ 相伴的对角线，则得到 $(x(1+y^{-1})^{-1},y^{-1})$。注意这些变换只是双有理的（例如，表达式 $x(1+y^{-1})^{-1}$ 在 $y=-1$ 处没有定义）。

上述图像对一般的 $n$ 都成立：$\C\PP^1$ 中 $n$ 个点的模空间 $\mathcal{M}_n$ 是一个代数簇，其坐标卡由 $n$ 边形的三角剖分编号，转移函数是特殊的双有理变换。这类变换律的普遍出现正是簇代数簇这一概念之源。

\section{三角剖分、箭图与格}

我们从有向纤毛曲面 $S$ 上的一个理想三角剖分 $\Delta$ 这一初始数据出发。这个三角剖分提供的核心组合数据是一个有限集（即 $\Delta$ 的弧集），连同「两条边属于同一个三角形」这一关系。
我们可以用一类特殊的图——箭图——来表示由 $\Delta$ 给出的全部这些信息。

\begin{definition}[5.1]
箭图（quiver）是指无 1-循环与 2-循环的有向图。更确切地说，它是一个三元组 $Q=(Q_0,Q_0^f,Q_1)$，其中 $Q_0$ 是有限集，$Q_0^f\subset Q_0$，而 $Q_1$ 是有限个有序对 $(v,v')\in Q_0^2$ 构成的集合，使得若 $(v,v')\in Q_1$，则 $v\ne v'$ 且 $(v',v)\notin Q_1$。每个 $(v,v')\in Q_1$ 带有一个重数 $m_{(v,v')}>0$。若 $(v',v)\in Q_1$，则令 $m_{(v,v')}=-m_{(v',v)}$；若 $(v,v')$ 与 $(v',v)$ 都不属于 $Q_1$，则令 $m_{(v,v')}=0$。$Q_0$ 的元素称为顶点，$Q_0^f$ 中的顶点称为冻结顶点，而 $(v,v')\in Q_1$ 称为从 $v$ 到 $v'$ 的箭。
\end{definition}

与纤毛曲面的三角剖分 $\Delta$ 相伴的箭图 $Q_\Delta=(Q_0,Q_0^f,Q_1)$ 如下得到：$Q_0$ 是 $\Delta$ 的边集，冻结顶点集 $Q_0^f$ 是外边集，并且对 $\Delta$ 的任意有向三角形 $\gamma,\gamma',\gamma''$，有从 $\gamma$ 到 $\gamma'$、从 $\gamma'$ 到 $\gamma''$ 以及从 $\gamma''$ 到 $\gamma$ 的箭（见图 5.2）。

% [原文 p.56]
\Fig[0.7\linewidth]{figures/fig_5_2.png}{5.2}{与三角剖分 $\Delta$ 相伴的箭图 $Q_\Delta$。}

给定箭图 $Q=(Q_0,Q_0^f,Q_1)$，考虑由顶点 $(e_v)_{v\in Q_0}$ 自由生成的 $\Z$-模 $\Lambda$。我们称 $\Lambda$ 为与 $Q$ 相伴的格。这个格自然带有一个由
\[
(e_v,e_{v'})=m_{(v,v')}
\]
定义的斜对称双线性形式 $(\cdot,\cdot)$。

由构造，格 $\Lambda$ 还带有一个典范基 $(e_v)_{v\in Q_0}$，而冻结顶点子集 $Q_0^f$ 诱导出这一基的子族 $(e_v)_{v\in Q_0^f}$。

\begin{definition}[5.2]
一个种子（seed）是一个四元组 $\mathsf{S}=(\Lambda,(\cdot,\cdot),(e_i)_{i\in I},J)$，其中 $\Lambda$ 是一个格（即有限生成的自由 $\Z$-模），$(\cdot,\cdot)$ 是 $\Lambda^2$ 上的斜对称双线性形式，$(e_i)_{i\in I}$ 是 $\Lambda$ 的一组基，而 $J$ 是 $I$ 的一个（可能空的）子集。
\end{definition}

\begin{remark}[5.3]
还存在一个略更一般的种子概念，即在箭图的每个顶点上再附加一个正整数（称为乘子，multiplier）作为数据，但就本书所讨论的主题而言，我们不需要这种程度的一般性。关于这一构造的更多细节可参见 [132, 260]。
\end{remark}

注意，可以由种子 $(\Lambda,(\cdot,\cdot),(e_i)_{i\in I},J)$ 反读出箭图的数据，方式如下：$Q_0=I$，$Q_0^f=J$，且对所有的 $v,v'\in I$，有一条从 $v$ 到 $v'$、重数为 $(e_v,e_{v'})$ 的箭。

\section{簇种子环面与坐标}

每个种子定义两个值得关注的分裂代数环面 $(\C^\times)^n$。事实上，每个秩为 $n$ 的格 $\Lambda$ 都定义一个复分裂环面
\[
\mathcal{X}_\Lambda:=\Hom(\Lambda,\C^\times)\cong(\C^\times)^n .
\]
赋值 $\Lambda\mapsto\mathcal{X}_\Lambda$ 诱导了从（有限秩）格范畴到分裂代数环面范畴的反变等价。

任何种子 $(\Lambda,(\cdot,\cdot),(e_i)_{i\in I},J)$ 都定义了簇环面 $\mathcal{X}_\Lambda$ 以及 $\mathcal{A}_\Lambda:=\mathcal{X}_{\Lambda^*}$，其中 $\Lambda^*$ 是 $\Lambda$ 的对偶格。我们将看到，反对称双线性形式 $(\cdot,\cdot)$ 本身就在 $\mathcal{X}_\Lambda$ 与 $\mathcal{A}_\Lambda$ 上诱导出有趣的结构：前者带有一个泊松括号（一个同时是导子的李括号），后者带有一个（可能退化的）闭 2-形式。此外，$(\cdot,\cdot)$ 定义一个映射 $\mathcal{A}_\Lambda\to\mathcal{X}_\Lambda$。

% [原文 p.57]
首先注意，任何 $\lambda\in\Lambda$ 都定义一个元素 $x_\lambda\in\Hom(\mathcal{X}_\Lambda,\C^\times)$，其定义为
\[
\begin{aligned}
x_\lambda:\ \mathcal{X}_\Lambda&\longrightarrow\C^\times\\
x&\longmapsto x(\lambda).
\end{aligned}
\]
对所有 $\lambda,\mu\in\Lambda$ 都有 $x_\lambda x_\mu=x_{\lambda+\mu}$。由种子给出的 $\Lambda$ 的基 $\{e_1,\dots,e_n\}$ 定义了 $\mathcal{X}_\Lambda$ 上的 $\C^\times$-值坐标，我们把它记作 $x_1,\dots,x_n$。此外，对所有的 $i,j=1,\dots,n$，记 $\epsilon_{ij}:=(e_i,e_j)$。

\begin{proposition}[5.4]
形式 $(\cdot,\cdot)$ 在 $\mathcal{X}_\Lambda$ 上诱导一个泊松括号，其定义为
\[
\{x_\lambda,x_\mu\}=(\lambda,\mu)\cdot x_\lambda x_\mu.
\]
\end{proposition}

\begin{proof}
容易验证 $\{x_\lambda x_{\lambda'},x_\mu\}=x_\lambda\{x_{\lambda'},x_\mu\}+x_{\lambda'}\{x_\lambda,x_\mu\}$，并且 $\{\cdot,\cdot\}$ 满足使其成为泊松括号所需的其余一切性质。在坐标下，$\{\cdot,\cdot\}$ 写作
\[
\{x_i,x_j\}=\epsilon_{ij}x_ix_j.
\]
\end{proof}

现在考虑对偶格 $\Lambda^*=\Hom(\Lambda,\Z)$，并令 $\{e_1^*,\dots,e_n^*\}$ 为 $\{e_1,\dots,e_n\}$ 在 $\Lambda^*$ 中的对偶基。与上面类似，这组基定义了
\[
\mathcal{A}_\Lambda:=\Hom(\Lambda^*,\C^\times)\cong(\C^\times)^n
\]
上的 $\C^\times$-值坐标 $a_1,\dots,a_n$。

\begin{proposition}[5.5]
$\Lambda$ 上的形式 $(\cdot,\cdot)$ 在 $\mathcal{A}_\Lambda$ 上诱导一个闭 2-形式。
\end{proposition}

\begin{proof}
在坐标下，形式 $(\cdot,\cdot)$ 写作
\[
(\cdot,\cdot)=\sum_{i<j}\epsilon_{ij}\,e_i^*\wedge e_j^*\in\bigwedge^{2}\Lambda^* .
\]
对偶地，它定义了 $\mathcal{A}_\Lambda$ 上的函数
\[
\sum_{i<j}\epsilon_{ij}\,a_i\wedge a_j .
\]
注意坐标 $a_i$ 在 $\mathcal{A}_\Lambda$ 上永不取零，因为它们是 $\C^\times$-值的。虽然每个 $a_i$ 的对数 $\log(a_i)$ 只定义到相差 $2\pi$ 的平移，但 1-形式 $\dd\log(a_i)$ 在 $\mathcal{A}_\Lambda$ 上无歧义地定义。因此
\[
\sum_{i<j}\epsilon_{ij}\,\dd\log(a_i)\wedge\dd\log(a_j)
\]
是 $\mathcal{A}_\Lambda$ 上的 2-形式，并且它显然是闭的。

% [原文 p.58]
事实上，这个 2-形式可以不依赖坐标地定义。双线性反对称形式 $(\cdot,\cdot)$ 诱导出元素 $W\in\OO(\mathcal{A}_\Lambda)^*\wedge\OO(\mathcal{A}_\Lambda)^*$，其中 $\OO(\mathcal{A}_\Lambda)^*$ 是 $\mathcal{A}_\Lambda$ 上可逆（正则）函数构成的群。对 $\OO(\mathcal{A}_\Lambda)^*\wedge\OO(\mathcal{A}_\Lambda)^*$ 的任意元素作用 $\dd\log\wedge\dd\log$ 都给出 $\mathcal{A}_\Lambda$ 上的一个 2-形式；对 $W$ 自然也是如此。更多细节可参见 [132]。
\end{proof}

% [Tutor Note: 命题 5.5 证明中「Dually, it defines the function $\sum_{i<j}\epsilon_{ij}a_i\wedge a_j$」一句须按双向量理解：$a_i$ 是 $\mathcal{A}_\Lambda$ 上的函数而非 1-形式，字面的 $\wedge$ 并不合法。正确表述是：$(\cdot,\cdot)$ 对偶地给出 $\mathcal{A}_\Lambda$ 上的双向量 $\sum_{i<j}\epsilon_{ij}a_ia_j\,\partial_{a_i}\wedge\partial_{a_j}$（即泊松结构），其经 $\dd\log\wedge\dd\log$ 或等价地在 $\dd\log a_i$ 上取值即得闭 2-形式 $\sum_{i<j}\epsilon_{ij}\dd\log a_i\wedge\dd\log a_j$。核对方法：$\dd(\dd\log a_i)=0$ 说明所得的 2-形式自动闭，与命题的断言一致。]
\begin{proposition}[5.6]
存在自然的映射 $p:\mathcal{A}_\Lambda\to\mathcal{X}_\Lambda$。
\end{proposition}

\begin{proof}
$\Lambda$ 上的双线性形式 $(\cdot,\cdot)$ 通过
\[
\lambda\longmapsto\sum_i(\lambda,e_i)\,e_i^*
\]
诱导映射 $\Lambda\to\Lambda^*$。
由于函子 $\Hom(\cdot,\C^\times)$ 是反变的，我们得到所需的 $p:\mathcal{A}_\Lambda\to\mathcal{X}_\Lambda$。在坐标下，映射 $p$ 写作
\[
p(a_1,\dots,a_n)=\Big(\prod_j a_j^{\epsilon_{1j}},\dots,\prod_j a_j^{\epsilon_{nj}}\Big).
\]
\end{proof}

\section{变异}

在本节中，我们关心前面引入的对象对所选三角剖分的依赖。确实，簇代数簇构造的目标是建立只依赖曲面、而不依赖理想三角剖分选取的对象。

首先回忆定理 4.20：纤毛曲面 $S$ 上任意两个理想三角剖分都可以通过有限次称为翻转的初等变换相联系。该定理意味着，只需理解一次三角剖分翻转所引起的改变，就能理解任意的三角剖分改变。

第一步是把翻转操作写成箭图的语言。于是我们现在定义箭图变异（mutation）的概念。设 $Q=(Q_0,Q_0^f,Q_1)$ 是箭图，$i\in Q_0\setminus Q_0^f$ 是一个非冻结顶点。在顶点 $i$ 处变异 $Q$ 所得的箭图 $Q'=\mu_i(Q)$ 构造如下：
\begin{enumerate}
\item 把所有以 $i$ 为起点或终点的箭反向（保持其重数）。
\item 对所有形如 $j\to i\to k$（$j,k\in Q_0$）的箭对，新增一条从 $k$ 到 $j$ 的箭，其重数 $m_{(k,j)}=m_{(j,i)}m_{(i,k)}$。
\item 若在这些修改之后，存在顶点对 $(j,k)$ 同时有一条重数为 $m_{(j,k)}$ 的箭 $j\to k$ 和一条重数为 $m_{(k,j)}$ 的箭 $k\to j$，且 $m_{(j,k)}<m_{(k,j)}$，则删去箭 $j\to k$，并把 $k\to j$ 的重数设为 $m_{(k,j)}-m_{(j,k)}$。若 $m_{(j,k)}=m_{(k,j)}$，则两条箭都被删去。最后这一步等价于计算箭的代数重数。
\end{enumerate}

% [原文 p.59]
\begin{example}[5.7]
让我们看一个例子：与四边形的理想三角剖分 $\Delta$ 相伴的箭图 $Q_\Delta$ 的一次变异。在任一非冻结顶点处的变异在图 5.3 中被逐步描述。注意变异之后所得的箭图，正是把被变异的那条边翻转后所得三角剖分相伴的箭图。同一现象对任意纤毛曲面的任意理想三角剖分都成立。
注意，虽然给定纤毛曲面 $S$ 通常有无穷多个不同的理想三角剖分，但从 $S$ 的三角剖分能得到的箭图只有有限多个。这是因为，如注 4.18 所述，理想三角剖分在纯映射类群作用下的轨道只有有限多个，而同一轨道中的箭图都同构。
\end{example}

\Fig[0.95\linewidth]{figures/fig_5_3.png}{5.3}{与四边形的一个三角剖分相伴的箭图的一次变异。变异所得箭图正是与四边形的另一个三角剖分相伴的箭图。}

为了计算方便，我们还需要知道变异对与箭图相伴的种子的作用。变异对种子 $\mathsf{S}=(\Lambda,(\cdot,\cdot),(e_i)_{i\in I},J)$ 的唯一改变发生在基 $e=(e_i)_{i\in I}$ 上。设 $i\in I$。由变异 $\mu_i$ 得到的新基 $\mu_i(e)$ 为
\[
\big(\mu_i(e)\big)_k=\begin{cases}
e_k+\max(\epsilon_{ki},0)e_i, & \text{若 } k\ne i\\
-e_i, & \text{若 } k=i .
\end{cases}
\]
我们记 $\mu_i(\mathsf{S})$ 为新种子 $(\Lambda,(\cdot,\cdot),\mu_i(e),J)$。注意，虽然斜对称双线性形式在变异下不变，但它在种子所给基下的矩阵 $\epsilon$ 确实会改变。$(\cdot,\cdot)$ 在基 $\mu_i(e)$ 下的矩阵 $\mu_i(\epsilon)$ 由下式给出：
\[
\big(\mu_i(\epsilon)\big)_{k\ell}=\begin{cases}
-\epsilon_{k\ell}, & \text{若 } i\in\{k,\ell\}\\
\epsilon_{k\ell}, & \text{若 } \epsilon_{ki}\epsilon_{i\ell}\le 0,\ i\notin\{k,\ell\}\\
\epsilon_{k\ell}+|\epsilon_{k,i}|\epsilon_{i\ell}, & \text{若 } \epsilon_{ki}\epsilon_{i\ell}>0,\ i\notin\{k,\ell\} .
\end{cases}
\]
由于种子的变异既不改变格也不改变斜对称双线性形式，两个环面 $\mathcal{X}_\Lambda$ 与 $\mathcal{A}_\Lambda$ 连同它们所带的结构都在变异下不变：

% [Tutor Note: 矩阵 $\mu_i(\epsilon)$ 第三行「$\epsilon_{k\ell}+|\epsilon_{k,i}|\epsilon_{i\ell}$」与 Fomin–Zelevinsky 的标准变异公式 $b'_{k\ell}=b_{k\ell}+\frac{1}{2}\big(|b_{ki}|b_{i\ell}+b_{ki}|b_{i\ell}|\big)$ 形式不同但等价：条件 $\epsilon_{ki}\epsilon_{i\ell}>0$ 保证二者同号，于是 $|b_{ki}|b_{i\ell}=b_{ki}|b_{i\ell}|=|\epsilon_{k,i}|\epsilon_{i\ell}$，两种写法给出同一个数。核对：取 $\epsilon_{ki}=\epsilon_{i\ell}=1$，标准式给 $b_{k\ell}+1$，原文式给 $\epsilon_{k\ell}+1$；取二者同为 $-1$，标准式给 $b_{k\ell}-1$，原文式 $|\epsilon_{k,i}|\epsilon_{i\ell}=-\epsilon_{ki}\epsilon_{i\ell}=-1$，亦相符。]
\begin{proposition}[5.8]
设 $\mathsf{S}$ 是一个种子。$\mathcal{X}_\Lambda$ 上的泊松括号与 $\mathcal{A}_\Lambda$ 上的闭 2-形式在种子 $\mathsf{S}$ 的变异下不变。
\end{proposition}

% [原文 p.60]
然而，$\mathcal{X}_\Lambda$ 上的 $\C^\times$-值坐标 $(x_i)_{i\in I}$ 与 $\mathcal{A}_\Lambda$ 上的 $(a_i)_{i\in I}$ 会按如下方式被变异改变：
\begin{equation}
\mu_i(x_j)=\begin{cases}
x_i^{-1}, & \text{若 } j=i\\
x_j\big(1+x_i^{\sgn(\epsilon_{ij})}\big)^{\epsilon_{ij}}, & \text{若 } j\ne i,
\end{cases}\tag{5.1}
\end{equation}
以及
\begin{equation}
\mu_i(a_j)=\begin{cases}
a_i^{-1}\Big(\displaystyle\prod_{k\mid \epsilon_{ki}>0}a_k^{\epsilon_{ki}}+\prod_{k\mid \epsilon_{ki}<0}a_k^{-\epsilon_{ki}}\Big), & \text{若 } j=i\\[3pt]
a_j, & \text{若 } j\ne i .
\end{cases}\tag{5.2}
\end{equation}
这两个公式分别称为 X 型与 A 型的变异公式。

\begin{remark}[5.9]
特别地，若 $x_1,\dots,x_n$（相应地 $a_1,\dots,a_n$）是点 $x\in\mathcal{X}_\Lambda$（相应地 $a\in\mathcal{A}_\Lambda$）在给定种子 $\mathsf{S}$ 下的坐标，则对由 $\mathsf{S}$ 经逐次变异得到的任意种子 $\mathsf{S}'$，$x$ 的坐标 $x_1',\dots,x_n'$（相应地 $a$ 的坐标 $a_1',\dots,a_n'$）都是 $x_1,\dots,x_n$（相应地 $a_1,\dots,a_n$）的有理函数。
\end{remark}

正如 4.3 节所见，把理想三角剖分 $\Delta$ 变成另一个 $\Delta'$ 的任意两个翻转序列只相差有限多个交换关系与五边形关系。由这些公式可知，对任意给定的三角剖分 $\Delta$ 及其相伴种子 $\mathsf{S}$，以及对 $\Delta$ 中任意两条不属于同一个三角形的弧 $i,j$，对所有 $k\in\Delta$ 都有 $\mu_i\circ\mu_j(x_k)=\mu_j\circ\mu_i(x_k)$ 与 $\mu_i\circ\mu_j(a_k)=\mu_j\circ\mu_i(a_k)$。

% [Tutor Note: 式 (5.1) 的两个分支与第 5.1 节的两个翻转公式可以互相核对。设五边形两条对角线对应的坐标是 $x,y$：当 $\epsilon_{xy}=+1$ 时 $\mu_x$ 给出 $(x^{-1},\,y(1+x^{\sgn(+1)})^{+1})=(x^{-1},y(1+x))$，当 $\epsilon_{xy}=-1$ 时 $\mu_x$ 给出 $(x^{-1},\,y(1+x^{-1})^{-1})$ —— 这正是例 5.13 列出的两个坐标变换。而原文 p.55（PDF p95）的第二个式子 $(x(1+y^{-1})^{-1},y^{-1})$ 则是变异 $y$ 所在对角线（即 $\mu_y$，取 $\epsilon_{yx}=-1$）的结果；把该式的 $x,y$ 互换即回到同一族变换。要点：$\sgn$ 与指数 $\epsilon_{ij}$ 的配合会在 $\epsilon_{ij}$ 取负时自动把 $1+x_i$ 换成 $1+x_i^{-1}$，不必另设公式。]
\begin{proposition}[{5.10（五边形关系）}]
设 $\Delta$ 是 $S$ 的理想三角剖分，$\mathsf{S}$ 是相伴的种子。设 $i,j\in\Delta$ 使得 $m_{i,j}=1$。则对所有 $k\in\Delta$ 有
\[
\mu_i\circ\mu_j\circ\mu_i\circ\mu_j\circ\mu_i(x_k)=\begin{cases}
x_k, & \text{若 } k\ne i,j\\
x_j, & \text{若 } k=i\\
x_i, & \text{若 } k=j,
\end{cases}
\]
以及
\[
\mu_i\circ\mu_j\circ\mu_i\circ\mu_j\circ\mu_i(a_k)=\begin{cases}
a_k, & \text{若 } k\ne i,j\\
a_j, & \text{若 } k=i\\
a_i, & \text{若 } k=j .
\end{cases}
\]
\end{proposition}

\begin{proof}
证明是直接的，但涉及相当长的计算。我们在此给出五边形上 A 型坐标的计算，其中所有外边的坐标都冻结且等于 $1$（见图 5.4）。于是两个内部 A 坐标按如下方式变化：
\begin{align*}
(a_i,a_j)&\xrightarrow{\ \mu_i\ }\Big(\frac{a_j+1}{a_i},a_j\Big)
\xrightarrow{\ \mu_j\ }\Big(\frac{a_j+1}{a_i},\frac{1+a_i+a_j}{a_ia_j}\Big)
\xrightarrow{\ \mu_i\ }\Big(\frac{a_i+1}{a_j},\frac{1+a_i+a_j}{a_ia_j}\Big)\\
&\xrightarrow{\ \mu_j\ }\Big(\frac{a_i+1}{a_j},a_i\Big)
\xrightarrow{\ \mu_i\ }(a_j,a_i).
\end{align*}
\end{proof}

% [原文 p.61]
\Fig[0.72\linewidth]{figures/fig_5_4.png}{5.4}{与五边形关系相伴的变异序列。注意这一变异序列使两条内边互换。}

变异公式的基本性质之一是：由于它们不涉及任何减法，它们保持实正坐标的集合：若 $\mathcal{X}_\Lambda$ 或 $\mathcal{A}_\Lambda$ 中的点在某个种子 $\mathsf{S}$ 下的所有坐标都是实正的，则它在由 $\mathsf{S}$ 经变异得到的任何种子下都将具有实正坐标。这一性质将在 10.2 节把簇代数簇与高阶泰希米勒理论联系起来时起根本作用。

A 型坐标的另一个重要性质是所谓的 Laurent 现象（Laurent phenomenon）：

\begin{theorem}[{5.11（见 {{{{{[133]}}}}}）}]
设 $\mathsf{S}$ 是种子，$a_1,\dots,a_n$ 是点 $a\in\mathcal{A}_\Lambda$ 在 $\mathsf{S}$ 下的坐标。设 $\mathsf{S}'$ 是由 $\mathsf{S}$ 经逐次变异得到的种子。则点 $a$ 的坐标 $a_1',\dots,a_n'$ 是变量 $a_1,\dots,a_n$ 的 Laurent 多项式。
\end{theorem}

A 坐标的这一性质对 X 坐标也有推论，尽管不那么直接，更多细节参见 [132]。

\section{簇代数簇}

\begin{definition}[5.12]
设 $\mathsf{S}$ 是种子，$|\mathsf{S}|$ 表示由 $\mathsf{S}$ 经一列变异能得到的所有种子的集合。定义 $\mathsf{S}$ 的 A 型簇代数簇（A-variety）为空间
% [原文 p.62]
\[
\mathcal{A}_{\mathsf{S}}:=\Big(\bigsqcup_{\mathsf{S}'\in|\mathsf{S}|}\mathcal{A}_{\Lambda_{\mathsf{S}'}}\Big)\Big/\sim,
\]
其中 $\Lambda_{\mathsf{S}'}$ 表示种子 $\mathsf{S}$ 的格，而 $a'\in\mathcal{A}_{\Lambda_{\mathsf{S}'}}\sim a''\in\mathcal{A}_{\Lambda_{\mathsf{S}''}}$ 当且仅当存在某个 $i$ 使得 $\mu_i(\mathsf{S}')=\mathsf{S}''$ 且 $\mu_i(a')=a''$。类似地，定义 $\mathsf{S}$ 的 X 型簇代数簇（X-variety）为
\[
\mathcal{X}_{\mathsf{S}}:=\Big(\bigsqcup_{\mathsf{S}'\in|\mathsf{S}|}\mathcal{X}_{\Lambda_{\mathsf{S}'}}\Big)\Big/\sim,
\]
其中 $x'\in\mathcal{X}_{\Lambda_{\mathsf{S}'}}\sim x''\in\mathcal{X}_{\Lambda_{\mathsf{S}''}}$ 当且仅当存在某个 $i$ 使得 $\mu_i(\mathsf{S}')=\mathsf{S}''$ 且 $\mu_i(x')=x''$。
\end{definition}

% [Tutor Note: 定义 5.12 中「$\Lambda_{\mathsf{S}'}$ 表示种子 $\mathsf{S}$ 的格」一句，按记号本身应为「种子 $\mathsf{S}'$ 的格」：$\Lambda_{\mathsf{S}'}$ 的下标是 $\mathsf{S}'$，且它在并集中随 $\mathsf{S}'\in|\mathsf{S}|$ 变化；若取 $\mathsf{S}$ 的格则与 $\mathsf{S}'$ 无关，并集就无从区分各个环面。此处照原文译出，读者按 $\Lambda_{\mathsf{S}'}$ 的下标理解即可。]
种子 $\mathsf{S}'\in|\mathsf{S}|$ 对应于 A 型簇与 X 型簇上一个取值于 $(\C^\times)^n$ 的开稠密坐标卡，它分别由与种子 $\mathsf{S}'$ 相伴的 A 坐标与 X 坐标给出。坐标卡之间的转换由变异函数给出，从而是双有理的。由于变异函数保持 A 环面上的闭 2-形式 $(\cdot,\cdot)$ 与 X 环面上的泊松括号 $\{\cdot,\cdot\}$，这些结构分别延拓为 $\mathcal{A}_{\mathsf{S}}$ 上的闭 2-形式 $(\cdot,\cdot)$ 与 $\mathcal{X}_{\mathsf{S}}$ 上的泊松括号 $\{\cdot,\cdot\}$。命题 5.6 中定义的映射 $p$ 也延拓为一个映射
\[
p:\mathcal{A}_{\mathsf{S}}\to\mathcal{X}_{\mathsf{S}} .
\]

\begin{example}[5.13]
让我们回到 5.1 节给出的 $\C^2$ 中 5 条直线模空间这一基本例子，并把它描述为 X 型簇代数簇的结构。
与 $\mathcal{M}_5$ 相伴的箭图是一个被三角剖分的五边形的箭图。它有 7 个顶点，其中 5 个是冻结的（即对应五边形外边的那几个）。这个箭图也给出初始种子 $(\Lambda,(\cdot,\cdot),(e_i)_{i\in Q_0},Q_0^f)$。
关键是：给出坐标变换 $(x,y)\mapsto(x^{-1},y(1+x))$ 或 $(x,y)\mapsto(x^{-1},y(1+x^{-1})^{-1})$ 的交比变换，恰好就是式 (5.1) 中的 X 型变异。因此，把与初始种子相伴的 X 环面粘到由一次翻转得到的另一个种子，正好表示两个坐标卡之间的过渡。
同样的方法可用于描述 $\C^2$ 中 $n$ 条直线的模空间，此时使用 $n$ 边形的三角剖分。簇理论表明该模空间是一个泊松簇。
\end{example}

\begin{example}[5.14]
与 $\C\PP^1$ 中 $n$ 个点的模空间非常类似的是 $\C^2$ 中 $n$ 个向量的模空间，它是一个 A 型簇代数簇。
给 $\C^2$ 配备一个复体积形式 $\omega\in\Omega^2(\C^2)$。我们考虑 $\C^2$ 中 $n$ 个向量 $(v_1,\dots,v_n)$ 构成的空间 $\mathcal{V}_n$ 模去 $\SL_n(\C)$ 的自然作用，使用 GIT 商。
我们再次考虑一个 $n$ 边形，把向量 $v_i$ 对应到顶点 $i$。固定多边形的一个三角剖分。对每条边（也包括多边形的外部边），我们配以 $\omega(v_i,v_j)\in\C$，其中 $v_i$ 与 $v_j$ 是该边两个顶点处的向量。但这里有一个明显的问题：由于 $\omega$ 是反对称的，$\omega(v_j,v_i)=-\omega(v_i,v_j)$，所以这个标号只定义到相差一个符号。为解决这个问题，我们实际上必须把多边形顶点 $i$ 与两个向量 $v_i$ 和 $v_{n+i}:=-v_i$ 相联系。
\end{example}

% [原文 p.63]
规则如下：取所选边的一个端点处的一个向量 $v_k\in\{v_i,v_{n+i},v_j,v_{n+j}\}$，取沿顺时针方向到另一端点所经的距离 $d$，并配以 $\omega(v_k,v_{k+d})$，其中 $k+d$ 按模 $2n$ 计算。这样，边上的坐标就有良好定义。

可以验证，这 $n$ 个向量可以由这些坐标重构出来，而且一次翻转将对坐标诱导一次 A 型簇变异，如式 (5.2) 所示，其中箭图与 $\C^2$ 中 $n$ 条直线模空间所用的箭图相同。因此，模空间 $\mathcal{V}_n$ 是一个 A 型簇代数簇。注意存在一个典范映射 $\pi:\mathcal{V}_n\to\mathcal{M}_n$，由取每个向量张成的直线得到。对应的 A 型簇与 X 型簇之间的典范映射 $p$ 满足：若 $a\in\mathcal{A}$ 是 $n$ 元向量组 $v=(v_1,\dots,v_n)\in\mathcal{V}_n$ 的坐标，则 $x=p(a)\in\mathcal{X}$ 是 $\pi(v)\in\mathcal{M}_n$ 的坐标。

% [Tutor Note: 例 5.14 中「模去 $\SL_n(\C)$ 的自然作用」按上下文应为 $\SL_2(\C)$：向量都在 $\C^2$ 中，体积形式取 $\omega\in\Omega^2(\C^2)$，只有 $\SL_2(\C)$ 在 $\C^2$ 上有自然作用并保持 $\omega$；$\SL_n(\C)$ 一般并不作用在 $\C^2$ 上（除非 $n=2$）。此处照原文译出。另可核对引入 $v_{n+i}:=-v_i$ 的必要性：因为 $\omega(v_{n+i},v_j)=-\omega(v_i,v_j)$，符号歧义正是由「在 $\{v_i,v_{n+i},v_j,v_{n+j}\}$ 中取 $v_k$、再按 $k+d\bmod 2n$ 走顺时针距离」这条规则消掉的。]
如前所述，由于变异公式不涉及减法，若把域 $\C$ 换成半域（semifield），例如 $\R_{>0}$，上述所有定义仍然有意义。这一特定的半域在研究高阶泰希米勒空间时将非常重要；见 10.2 节。

具有几何意义的其他半域是热带半域（tropical semifield），例如 $\Z^{\mathrm{t}}=(\Z,\max,+)$。在簇经典泰希米勒空间这一具体情形中，这些簇代数簇的热带点体现了曲面上测度层（measured lamination）的概念 [131]。这些对象出现在泰希米勒空间的 Thurston 紧化的边界上。用热带对象紧化模空间是一种常见做法。对高阶泰希米勒空间，热带版本在代数上有良好定义，尽管其几何解释只被部分理解 [371, 243, 130]。

\begin{remark}[5.15]
热带变异公式在超对称量子场论中有自然的物理实现，例如作为 $4d$ $\mathcal{N}=1$ 箭图规范理论中的 Seiberg 对偶；这一主题我们将在第 17 章讨论。
\end{remark}

对应于变异类 $|\mathsf{S}|$ 的簇代数簇 $\mathcal{A}_{\mathsf{S}}$ 与 $\mathcal{X}_{\mathsf{S}}$ 之间由一个称为簇对偶（cluster duality）的深刻对偶概念相联系，它在 [129, 猜想 12.3] 中被猜想，并在 [174] 中被（部分地）证明。泰希米勒空间的具体情形在 [131] 中处理。每个簇代数簇都带有一个特殊的函数空间：在每个簇坐标卡中都是系数为正的 Laurent 多项式的正则函数。这个空间是一个锥。换言之，簇对偶断言：对簇 $\mathcal{A}_{\mathsf{S}}$（相应地 $\mathcal{X}_{\mathsf{S}}$）而言，这一锥的极值元素与簇 $\mathcal{X}_{\mathsf{S}}$（相应地 $\mathcal{A}_{\mathsf{S}}$）的 $\Z^{\mathrm{t}}$-热带点一一对应。

对偶猜想给出的特殊基有一个显著性质：若 $(f_\mu)$ 表示这样一组基，则由
\[
f_\mu f_\nu=\sum_\rho c^\rho_{\mu,\nu}f_\rho
\]
得到的结构常数 $c^\rho_{\mu,\nu}$ 是正整数。这些结构常数的组合解释问题在一般情况下仍未解决。类似的现象出现在量子群的 Lusztig 典范基、Hecke 代数的 Kazhdan–Lusztig 基以及李群表示论中的融合系数里。

% [原文 p.64]
最后，我们想强调：簇代数簇的基本性质——即簇 A 型（相应地 X 型）簇代数簇带有闭 2-形式（相应地泊松括号）——是一些更深层结构的影子；关于这一点的一般讨论我们参见 [129]。具体而言，簇 A 型簇代数簇上的闭 2-形式实际上来自一个所谓的 motivic 化身（motivic avatar），而簇 X 型簇代数簇上的泊松结构容许一个典范量子化。在泰希米勒空间的语境中，这一量子化已在 [228, 132] 中明确实现。量子双对数
\[
\Psi_q(z)=\prod_{k=0}^{\infty}\frac{1}{1+q^{2k+1}z},
\]
起着核心作用，它既出现在量子变异规则中，也出现在映射类群对形变模空间的作用中。这些发展引出了量子 X 型簇代数簇（quantum X-variety）的概念：其中的变异规则被替换为用量子双对数作共轭。当 $q\to1$ 时的极限恢复经典的变异。

% [Tutor Note: 「$q\to1$ 的极限恢复经典变异」应理解为量子变异算子（用量子双对数作共轭）的经典极限，而不能逐字取 $\Psi_q$ 的极限：对 $|q|<1$，$\Psi_q(z)=\prod_{k\ge0}(1+q^{2k+1}z)^{-1}$ 收敛；但 $q\to1$ 时每个因子趋于 $(1+z)^{-1}$，无穷乘积发散。标准做法是令 $q=\ee^{\hbar}$ 并把 $\Psi_q$ 当作形式幂级数或算子（Faddeev 量子双对数的形式版本），经典极限对应 $\hbar\to0$；此时共轭作用对单项式给出式 (5.1)、(5.2) 那样的有理变换。]
